% Part: set-theory
% Chapter: ordinals
% Section: vn

\documentclass[../../../include/open-logic-section]{subfiles}

\begin{document}

\olfileid{sth}{ordinals}{vn}
\olsection[ফন নয়মানের নির্মাণ]{ফন নয়মানের ক্রমসংখ্যা-নির্মাণ}

\olref[sth][ordinals][iso]{thm:woalwayscomparable} থেকে
একটি ভাবনা আসে। সুক্রমবিন্যাসগুলির প্রতিনিধি হিসেবে
কিছু বস্তু আনতে পারি, যাদের বলব \emph{ক্রমপ্রকার}।
সুক্রমবিন্যাস $\tuple{A, <}$-এর ক্রমপ্রকার
$\ordtype{A, <}$ লিখলে, আমরা নিচের দুটি নীতি চাইব:
\begin{align*}
	\ordtype{A, <} = \ordtype{B, \lessdot} &
	\text{ যদি এবং কেবল যদি } \ordeq{\tuple{A, <}}{\tuple{B, \lessdot}}\\
	\ordtype{A, <} < \ordtype{B, \lessdot}&
	\text{ যদি এবং কেবল যদি }\ordeq{\tuple{A, <}}{\tuple{B_b, \lessdot_b}}\text{ কোনো }b \in B\text{-এর জন্য}
\end{align*}
তাছাড়া, স্বাভাবিক সংখ্যাগুলিকে যেমন নির্দিষ্ট কিছু
সেট হিসেবে আনা যায়, তেমনি ক্রমপ্রকারগুলিকেও
\emph{নির্দিষ্ট কিছু সেট} হিসেবে আনতে চাইতে পারি।

এটি করার সবচেয়ে প্রচলিত পদ্ধতি—এবং এখানে আমরা
সেই পথেই চলব—হলো কয়েকটি \emph{প্রামাণিক}
সুক্রমিত সেট দিয়ে ক্রমপ্রকারগুলি সংজ্ঞায়িত করা।
এই প্রামাণিক সেটগুলি প্রথম ফন নয়মান প্রবর্তন করেন:

\begin{defn}
$A$ সেটটি \emph{পরিযায়ী} {যদি এবং কেবল যদি}
$(\forall x \in A)x \subseteq A$। এরপর $A$ একটি
\emph{ক্রমসংখ্যা} {যদি এবং কেবল যদি} $A$ পরিযায়ী
এবং $\in$ দ্বারা সুক্রমিত।
\end{defn}
\noindent
এর পর ক্রমসংখ্যার জন্য গ্রিক অক্ষর ব্যবহার করব।
সংজ্ঞা থেকে সঙ্গে সঙ্গেই পাই: $\alpha$ যদি ক্রমসংখ্যা
হয়, তবে $\tuple{\alpha, \in_\alpha}$ একটি
সুক্রমবিন্যাস, যেখানে $\in_\alpha =
\Setabs{\tuple{x, y} \in \alpha^2}{x \in y}$।
তাই সংকেতলিপিতে একটু শিথিলতা মেনে $\alpha$
\emph{নিজেকেই} সুক্রমবিন্যাস বলতে পারি।

প্রথম কয়েকটি ক্রমসংখ্যা এই:
\[
	\emptyset, \{\emptyset\},
	\{\emptyset, \{\emptyset\}\},
	\{\emptyset, \{\emptyset\}, \{\emptyset, \{\emptyset\}\}\}, \ldots
\]
দেখবে, অসীমতার স্বতঃসিদ্ধে অর্থাৎ $\omega$-র
ফন নয়মানীয় সংজ্ঞায় প্রথম যে কয়েকটি ক্রমসংখ্যা
পেয়েছিলাম, এগুলিই সেগুলি
(\olref[sth][z][infinity-again]{sec} দেখো)।
এটি কাকতালীয় নয়। ফন নয়মানের সংজ্ঞায়
স্বাভাবিক সংখ্যাগুলি ক্রমসংখ্যা, আবার
ট্রান্সফাইনাইট ক্রমসংখ্যারও জায়গা আছে।

বরাবরের মতো এবারও প্রশ্ন করা যায়: এগুলিই কি
\emph{প্রকৃত} ক্রমসংখ্যা? নাকি ফন নয়মান আমাদের
কিছু সেট দিয়েছেন, যেগুলিকে ক্রমসংখ্যা
\emph{হিসেবে} ধরতে পারি? এ নিয়ে আলোচনা
\olref[sfr][rel][ref]{sec},
\olref[sfr][arith][ref]{sec},
\olref[sfr][infinite][dedekindsproof]{sec} এবং
\olref[sth][z][nat]{sec}-এর আলোচনার মতোই;
তাই আবার বিস্তারিত বলব না। এর পর ``ক্রমসংখ্যা''
বলতে আমরা শুধু ``ফন নয়মানীয় ক্রমসংখ্যা''ই
বোঝাব।

\end{document}
